% !PRINTED 594
\weberpage{594}
\begin{center}
\textbf{\S\,160.\qquad \gesperrt{Prime ideals of the first degree in $\mathfrak{T}_{\mathfrak{m}}$.}}
\end{center}

What now matters above all is to seek out the primes $p$ that are decomposable in the field $\mathfrak{L}$ into prime ideals of the first degree. By § 122 and § 156 these are the primes $p$ that are representable by the principal form of the discriminant $D = 4\Delta$ or $16\Delta$. These are of the form:
\[
\begin{aligned}
1) \quad & p = y^2 - \Delta x^2, \qquad \Delta \equiv 0 \ (\mathrm{mod}\ 4) \text{ or } \equiv 1 \ (\mathrm{mod}\ 8),\\
2) \quad & p = y^2 - 4\Delta x^2, \qquad \Delta \equiv 5 \ (\mathrm{mod}\ 8).
\end{aligned}
\]
If $\Delta \equiv 1 \ (\mathrm{mod}\ 8)$, then the third coefficient of the form $(a, b, c)$ must be even. If $\Delta \equiv 0 \ (\mathrm{mod}\ 4)$, then we may assume $c$ even, by passing if necessary to the parallel form $(a, b+2a, a+b+c)$. If therefore $p$ splits into the two prime factors $\pi, \pi'$, then
\weq{1}{\begin{aligned}
1) \quad & \pi = y + x\sqrt{\Delta},\\
2) \quad & \pi = y + 2x\sqrt{\Delta},
\end{aligned}}
and in order to apply the multiplication formula § 157, (10), (17) to $\mu = \pi$, we have in the latter case to replace $x$ by $2x$. We may therefore in all cases assume $cx$ even, and hence $\varepsilon = \pm 1$.

The conditions for the modulus 2 that lie in the representations (1), 1), 2) can be summarised in the single one:
\[
\pi \equiv 1 \ (\mathrm{mod}\ 2).
\]
If we set, when $N(\pi) = p$, for the moment by way of abbreviation:
\weq{2}{\begin{aligned}
\varphi(x) &= A_1x + A_3x^3 + \cdots + A_{p-2}x^{p-2} + x^p,\\
\psi(x) &= 1 + A_{p-2}x^2 + \cdots + A_3x^{p-3} + A_1x^{p-1}, % source plus has a small damage tick; no mathematical prime is printed here
\end{aligned}}
then by § 157, (18), (19)
\weq{3}{A_1 = \pm\pi,}
and
\weq{4}{\pm \boldsymbol{s}(\pi u) = \frac{\varphi(\boldsymbol{s})}{\psi(\boldsymbol{s})},}
and from this there follows by differentiation with respect to $u$:
\weq{5}{\pm \frac{\pi \boldsymbol{s}'(\pi u)}{\boldsymbol{s}'(u)} = \frac{\varphi'(\boldsymbol{s})\psi(\boldsymbol{s}) - \varphi(\boldsymbol{s})\psi'(\boldsymbol{s})}{\psi(\boldsymbol{s})^2}.}
But by § 155, (5)
\[
\boldsymbol{s}'(u)^2 = 1 - \tfrac{1}{4}\lambda \boldsymbol{s}^2 + \boldsymbol{s}^4,
\]
%CONT>>
